\documentclass[11pt,a4paper]{article}
\usepackage{amsmath,amssymb,amsthm}
\usepackage{graphicx}
\usepackage[margin=2.2cm]{geometry}
\usepackage{xcolor}
\usepackage{titlesec}
\usepackage{fancyhdr}
\usepackage{booktabs}
\usepackage{enumitem}
\usepackage{float}
\usepackage{hyperref}
\hypersetup{colorlinks=true,linkcolor=blue!70!black,urlcolor=blue!70!black}

% Allow LaTeX more freedom in placing floats
\renewcommand{\topfraction}{0.9}
\renewcommand{\bottomfraction}{0.9}
\renewcommand{\textfraction}{0.07}
\renewcommand{\floatpagefraction}{0.7}
\setcounter{topnumber}{3}
\setcounter{bottomnumber}{3}
\setcounter{totalnumber}{5}

\definecolor{primary}{RGB}{0,51,102}
\definecolor{accent}{RGB}{204,102,0}

\titleformat{\section}{\Large\bfseries\color{primary}}{\thesection}{1em}{}
\titleformat{\subsection}{\large\bfseries\color{primary!80!black}}{\thesubsection}{1em}{}

\pagestyle{fancy}
\fancyhf{}
\fancyhead[L]{\small\color{primary}Essential Summary}
\fancyhead[R]{\small\color{primary}Li et al.\ (2026)}
\fancyfoot[C]{\thepage}
\renewcommand{\headrulewidth}{0.5pt}

\begin{document}

\begin{center}
{\LARGE\bfseries\color{primary}Essential Summary}\\[12pt]
{\large\bfseries Quantitative Analysis of Dynamical Bifurcations\\in a Coupled Smooth and Discontinuous Oscillator\\with High-order Nonlinear Damping}\\[8pt]
{\normalsize Zhenbo Li, Linxia Hou, Ruyue Peng}\\[4pt]
{\small\itshape Physica Scripta, Vol.\ 101, No.\ 12, 125205 (2026)\quad DOI: 10.1088/1402-4896/ae5134}\\[6pt]
{\small\rule{0.7\textwidth}{0.5pt}}\\[4pt]
{\small\color{accent}This is an author-prepared essential summary, not the original paper.\\Please cite the published version.}
\end{center}

\vspace{0.5em}

\section{Background and Motivation}

The coupled smooth and discontinuous (SD) oscillator, introduced by Han et al.\ (2012), consists of two SD oscillators coupled rigidly. Its restoring force contains \textbf{two irrational nonlinear terms} arising from the geometric configuration of oblique springs. These irrational terms make the system significantly harder to analyze than classical nonlinear oscillators.

When perturbed by high-order nonlinear damping of the form $(\mu_c + \mu_1 x + \mu_2 x^2 + \mu_3 x^3 + \mu_4 x^4)\dot{x}$, the oscillator exhibits rich dynamical behavior: limit cycles of varying stability, homoclinic and heteroclinic bifurcations, and multi-stable configurations (single-well, double-well, and triple-well potentials).

Classical quantitative methods---the elliptic function L--P method, the elliptic perturbation method, and the standard GHFP method---\textbf{cannot handle} this system because the two irrational terms render the key integrals analytically intractable. This work overcomes that barrier by extending the \textbf{modified generalized harmonic function perturbation (MGHFP)} method to this most challenging case.

\section{System Model}

The investigated oscillator is governed by:
\begin{equation}
\ddot{x} + (x+\beta)\frac{1}{\sqrt{(x+\beta)^2+\alpha^2}} + (x-\beta)\frac{1}{\sqrt{(x-\beta)^2+\alpha^2}} = \varepsilon(\mu_c + \mu_1 x + \mu_2 x^2 + \mu_3 x^3 + \mu_4 x^4)\dot{x}
\label{eq:model}
\end{equation}
where $\alpha$ is the smoothing parameter ($\alpha \to 0$ gives the non-smooth/discontinuous limit), $\beta$ is the coupling parameter controlling the potential well structure, $\varepsilon$ is a small positive perturbation parameter, $\mu_c$ is the control parameter, and $\mu_1$--$\mu_4$ are nonlinear damping coefficients.

\subsection{Triple-Well Potential Structure}

As $\alpha$ and $\beta$ vary, the oscillator's potential energy landscape undergoes qualitative changes. For the triple-well configuration, the system has \textbf{five equilibrium points}: two centers, two saddles, and an additional center---creating a landscape where both homoclinic and heteroclinic orbits can coexist at the same saddle point.

\begin{figure}[htbp]
\centering
\includegraphics[width=0.85\textwidth,height=0.35\textheight,keepaspectratio]{fig5-potential.png}
\caption{\textbf{Figure 5.} Potential energy diagram of the oscillator in the triple-well configuration. The five equilibrium points create a complex landscape with three potential wells separated by saddle points, enabling coexistence of homoclinic and heteroclinic orbits.}
\end{figure}

\section{The MGHFP Method --- Key Steps}

\subsection{Step 1: Nonlinear Time Transformation}

Introduce $\frac{d\varphi}{dt} = \Phi(\varphi)$ with $\varphi(0)=0$, $\varphi(t+\tau)=\varphi(t)+\pi$. The undamped system transforms into:
\begin{equation}
\frac{d}{d\varphi}(\Phi x') + g(x) = 0, \qquad x(\varphi) = a\cos^2\Phi(\varphi) + b
\end{equation}
where $g(x) = (x+\beta)\left[1 - \frac{\alpha^2}{(x+\beta)^2+\alpha^2}\right] + (x-\beta)\left[1 - \frac{\alpha^2}{(x-\beta)^2+\alpha^2}\right]$.

The function $\Phi_0(\varphi)$ is expanded as a Fourier series:
\begin{equation}
\Phi_0(\varphi) = \sum_{i=0}^{m} \left(p_{2i}\cos 2i\varphi + q_{2i}\sin 2i\varphi\right)
\end{equation}

\subsection{Step 2: Perturbation Expansion}

The amplitude and frequency are expanded in powers of $\varepsilon$:
\begin{equation}
a = a_0 + \varepsilon a_1 + \cdots, \qquad \Phi = \Phi_0 + \varepsilon\Phi_1 + \cdots
\end{equation}
Substituting into the damped system and collecting terms at each order of $\varepsilon$ yields a hierarchy of equations. The zeroth-order gives the unperturbed solution; the first-order determines $a_1$ and $\Phi_1$.

\subsection{Step 3: Composite Simpson Integration (Key Innovation)}

The Fourier coefficients $p_{2i}$ involve integrals of the irrational nonlinear terms that \textbf{cannot be evaluated analytically}. The MGHFP method's key innovation is the introduction of the \textbf{composite Simpson quadrature formula} to discretize these integrals into sums over equally-spaced nodes. The integration interval $[0, \pi]$ is divided into eight equal parts, yielding explicit expressions for each $p_{2i}$ in terms of $a_0$, $b$, $\alpha$, and $\beta$.

\section{Key Equations}

\subsection{Amplitude--Parameter Relation (Central Result)}

Setting the Fourier truncation order $m=4$, the analytical relationship between the limit cycle amplitude $A$ and the control parameter $\mu_c$ is:
\begin{equation}
\boxed{\mu_c = -\frac{4}{a_0^2(2p_0 - p_4)}\Big(E_1\mu_1 + E_2\mu_2 + E_3\mu_3 + E_4\mu_4\Big)}
\label{eq:amplitude}
\end{equation}
where $p_{2i}$ are Fourier coefficients computed via composite Simpson integration, and $E_1$--$E_4$ are explicit functions of $a_0$, $b$, and $p_{2i}$. For symmetric limit cycles: $a_0 = 2A$, $b = -A$. For asymmetric limit cycles: $a_0 = A - C$, $b = C$, where $C$ depends on $\alpha$ and $\beta$.

\subsection{Stability Criterion}

The stability of each limit cycle is determined by the characteristic quantity:
\begin{equation}
\boxed{h_0 = \frac{1}{\pi}\int_0^\pi \left[\mu_c + \sum_{i=1}^{4}\mu_i(a_0\cos^2\varphi + b)^i\right]d\varphi}
\label{eq:stability}
\end{equation}
Based on the qualitative theory of ODEs:
\begin{itemize}[leftmargin=2em]
\item $h_0 > 0$: the limit cycle is \textbf{unstable}
\item $h_0 = 0$: the limit cycle is \textbf{semi-stable} (bifurcation point)
\item $h_0 < 0$: the limit cycle is \textbf{stable}
\end{itemize}

\subsection{Homo-heteroclinic Bifurcation Thresholds}

Setting $b = h$ (saddle point ordinate) in Eq.~\eqref{eq:amplitude} yields $\mu_c^{\mathrm{hom}}$ for \textbf{homoclinic} bifurcation. Setting $b = h_2$ and $a_0 + b = h_1$ (connecting two saddles) yields $\mu_c^{\mathrm{het}}$ for \textbf{heteroclinic} bifurcation.

\subsection{First-Order Approximate Solution}

\begin{equation}
x = 2A\cos^2\Phi - A + \mathcal{O}(\varepsilon^2), \qquad \dot{x} = -2a_0[\Phi_0 + \varepsilon\Phi_1]\cos\varphi\sin\varphi + \mathcal{O}(\varepsilon^2)
\end{equation}

\section{Example 1: Single-Well Potential ($\alpha=2$, $\beta=0.5$)}

With $\mu_1=1$, $\mu_2=1$, $\mu_3=1$, $\mu_4=-1$, the amplitude--parameter curve reveals the complete limit cycle lifecycle:

\begin{figure}[htbp]
\centering
\includegraphics[width=\textwidth,height=0.38\textheight,keepaspectratio]{fig1-amplitude.png}
\caption{\textbf{Figure 1.} (Upper) Relationship between limit cycle amplitude $A$ and control parameter $\mu_c$. (Lower) Characteristic quantity $h_0$ determining stability. Key bifurcation points: $\mu^{(0)}=-0.1249$ (semi-stable cycle emerges), $\mu^{(1)}=-0.05$ (splits into stable + unstable), $\mu^{(2)}=0.02$ (unstable cycle collapses to singular point).}
\end{figure}

The lifecycle unfolds as follows:
\begin{itemize}[leftmargin=2em]
\item $\mu_c < \mu^{(0)} = -0.1249$: \textbf{no limit cycle exists}.
\item At $\mu_c = \mu^{(0)}$: a \textbf{semi-stable} limit cycle emerges (outer-stable, inner-unstable).
\item $\mu^{(0)} < \mu_c < 0$: the semi-stable cycle \textbf{splits into two}---a smaller unstable cycle ($h_0 > 0$) and a larger stable cycle ($h_0 < 0$).
\item At $\mu_c = 0$: the unstable cycle \textbf{collapses} to the singular point.
\item $\mu_c > 0$: only a single \textbf{stable} limit cycle persists, growing with $\mu_c$.
\end{itemize}

\begin{figure}[htbp]
\centering
\includegraphics[width=\textwidth,height=0.38\textheight,keepaspectratio]{fig2-limitcycles.png}
\caption{\textbf{Figure 2.} Phase portraits of limit cycles at different $\mu_c$ values ($\varepsilon=0.2$). Solid lines: Runge--Kutta numerical integration. Dotted lines: present MGHFP analytical method. The excellent agreement validates accuracy even at moderately large $\varepsilon$.}
\end{figure}

\section{Example 2: Triple-Well Potential}

The triple-well case is considerably more complex: five equilibrium points, coexistence of homoclinic and heteroclinic orbits, and both small and large limit cycles.

\subsection{Amplitude--Parameter Relation for Triple-Well}

\begin{figure}[htbp]
\centering
\includegraphics[width=\textwidth,height=0.38\textheight,keepaspectratio]{fig6-triplewell-amp.png}
\caption{\textbf{Figure 6.} Relationship between the amplitude $A$ and control parameter $\mu_c$ (upper), and the characteristic quantity $h_0$ (lower), for the triple-well oscillator. The evolution of both small limit cycles (within homoclinic orbits) and large limit cycles (enclosing all equilibria) is quantitatively captured.}
\end{figure}

\subsection{Limit Cycle Phase Portraits}

\begin{figure}[htbp]
\centering
\includegraphics[width=\textwidth,height=0.38\textheight,keepaspectratio]{fig7-triplewell.png}
\caption{\textbf{Figure 7.} Limit cycles for the triple-well oscillator at different $\mu_c$ values ($\varepsilon=1$). Solid lines: Runge--Kutta method. Dotted lines: MGHFP method. Despite the large perturbation parameter, the analytical results closely match the numerical reference.}
\end{figure}

\subsection{Homoclinic Orbits}

The homoclinic orbits connect a saddle point to itself. The MGHFP method provides analytical approximate solutions for these orbits, which are validated against numerical integration.

\begin{figure}[htbp]
\centering
\includegraphics[width=0.82\textwidth,height=0.40\textheight,keepaspectratio]{fig10-homoclinic.png}
\caption{\textbf{Figure 10.} Phase portraits of homoclinic solutions with different parameters ($\varepsilon=0.5$). Solid lines: Runge--Kutta. Dotted: MGHFP. The method accurately captures the homoclinic orbit geometry even at relatively large perturbation.}
\end{figure}

\subsection{Heteroclinic Orbits}

The heteroclinic orbits connect two distinct saddle points. For the triple-well system, these coexist with homoclinic orbits at the same saddle---a uniquely challenging scenario that the MGHFP method successfully handles.

\begin{figure}[htbp]
\centering
\includegraphics[width=0.82\textwidth,height=0.40\textheight,keepaspectratio]{fig11-heteroclinic.png}
\caption{\textbf{Figure 11.} Phase portraits of heteroclinic solutions with different parameters ($\varepsilon=0.1$). Solid lines: Runge--Kutta. Dotted: MGHFP. The analytical heteroclinic orbits connecting two saddle points are in excellent agreement with numerical results.}
\end{figure}

\section{Key Findings}

\begin{enumerate}[leftmargin=2em]
\item \textbf{Full quantitative prediction:} The MGHFP method achieves, for the first time, analytical quantitative analysis of limit cycles in a coupled SD oscillator with two irrational nonlinearities. Given system parameters, one can directly compute the number, amplitude, and existence interval of every limit cycle---no repetitive numerical search needed.

\item \textbf{Analytical stability criterion:} The characteristic quantity $h_0$ distinguishes stable, unstable, and semi-stable limit cycles analytically. Unstable and semi-stable limit cycles---exceedingly difficult to capture numerically---are reliably identified by this method.

\item \textbf{Simultaneous homo-heteroclinic bifurcation prediction:} For the triple-well system, homoclinic and heteroclinic orbits coexist at the same saddle. The MGHFP method predicts both bifurcation parameters simultaneously and provides analytical orbit solutions (Figures 10--11), a capability not available from prior approaches.

\item \textbf{Complete limit cycle lifecycle:} The entire lifecycle is quantitatively tracked---from birth (semi-stable bifurcation at $\mu^{(0)}$), through evolution (splitting into stable + unstable pair), to annihilation (collapse to singular point at $\mu^{(2)}$), as clearly shown in Figure 1.

\item \textbf{High-order damping applicability:} With quartic nonlinear damping, the method maintains reliable accuracy when $|\varepsilon\mu_i| \leq 0.5$. Beyond this range, accuracy is recovered by increasing the Fourier truncation order.
\end{enumerate}

\section{Method Limitations}

\begin{itemize}[leftmargin=2em]
\item \textit{Large perturbation parameter} ($\varepsilon > 0.5$): Relative error exceeds 5\%. Higher-order perturbation expansions can mitigate this.
\item \textit{Limit cycles crossing saddle points}: Velocity drops sharply near the saddle, causing abrupt waveform features. Higher Fourier truncation orders are needed.
\item \textit{Non-smooth potential wells} ($\alpha = 0$): Non-smoothness affects accuracy. A smaller $\varepsilon$ or higher truncation order is recommended.
\end{itemize}

\section{Applications}

The model---combining irrational nonlinear stiffness, high-order nonlinear damping, and multi-stable dynamics---has direct relevance to:

\textbf{(1) QZS vibration isolators:} Quasi-zero stiffness isolators with nonlinear electromagnetic damping share the irrational stiffness and high-order damping structure. The MGHFP method provides a quantitative tool for predicting their limit cycle behavior and bifurcation thresholds.

\textbf{(2) MEMS resonators:} In micro-electro-mechanical systems, nonlinear damping and multi-stable dynamics are increasingly exploited. The bifurcation analysis framework directly supports design and optimization where saddle-node bifurcations and multi-stability play critical roles.

\vfill
\begin{center}
\small\rule{0.5\textwidth}{0.3pt}\\[4pt]
\small\color{primary} Author: Zhenbo Li (lizhenbo@usc.edu.cn)\\
\small University of South China, School of Mathematics and Physics\\
\small Hunan Key Laboratory of Mathematical Modeling and Scientific Computing
\end{center}

\end{document}